\documentclass[11pt]{article}

% Use Helvetica font.
%\usepackage{helvet}

%\usepackage[T1]{fontenc}
%\usepackage[sc]{mathpazo}
\usepackage{color}
\usepackage{amsmath,amsthm,amssymb,multirow,paralist}
%\renewcommand{\familydefault}{\sfdefault}

% Use 1/2-inch margins.
\usepackage[margin=0.8in]{geometry}
\usepackage{hyperref}

\begin{document}

\begin{center}
{\Large \textbf{Midterm project}}

\vspace{10pt}

Nikhil Kumar

\vspace{10pt}

\textit{March 2026}
\end{center}

\section{Introduction}

The Remaining Useful Life (RUL) of the mechanical systems is a key issue,
which needs to be predicted in the context of predictive maintenance. Unex-
pected engine failures may cause a great deal of financial losses and even serious
safety risks in such industries as aviation. Conventional maintenance plans such
as planned and reactive maintenance are not usually efficient since they either
prematurely replace parts or maintain them without stopping failures. Conse-
quently, more and more interest in data-driven methods that can approximate
system health and anticipate failures before they happen is growing.

The project at hand is concerned with the prediction of Remaining Use-
ful Life of jet engines based on NASA CMAPSS (Commercial Modular Aero-
Propulsion System Simulation) data [1]. This data has multivariate time-series
data of various engines at different conditions up to failure. At each time step,
operational settings and several sensor measurements are available and the goal
is to make a forecast of the number of cycles to go before an engine fails.

This is a regression time-series problem, and the model has to acquire degra-
dation tendencies over time. To find a solution to this, we apply baseline ma-
chine learning models and a deep learning methodology. In particular, we com-
pare linear regression and random forest regression to a 1D Convolutional Neural
Network (CNN) that is implemented to extract the temporal relationships in
sequential data. This project aims at assessing the performance of these mod-
els and find out whether deep learning techniques can enhance the accuracy of
prediction when estimating the RUL.

\section{Related Work}

The prediction of Remaining Useful Life (RUL) is a predictive maintenance
technique that has been extensively studied and referenced in relation to its
importance in ensuring that systems are made more reliable and that operation
costs are minimized. Degradation is a characteristic of complex systems like
aircraft engines that is affected by changes in different sensor measurements;
therefore, degradation is more applicable to data-driven approaches.

As a basic approach to RUL prediction, traditional machine learning ap-
proaches have been used, which include linear regression and tree-based ap-
proaches. However, these approaches are deemed ineffective in capturing time-
based dependencies that are characteristic of sequence-based data. In addition,
to address this limitation of traditional approaches, recent studies have been
conducted on deep learning approaches, particularly those that are based on
sequence-based approaches like RNN and LSTM [2].

More recently, time-series prediction tasks (and RUL estimation) have also
been performed by using Convolutional Neural Networks (CNNs) [3]. CNNs
can capture local temporal features and have few parameters than recurrent
ones which is computationally efficient and yet capable of obtaining good per-
formance.

In the given project, we feed a 1D CNN model with the engine sensor data to
identify time-related trends in the data and compare its functionality to other
baseline models, including linear regression, and random forest regression. This
will enable us to assess the claim that deep learning methods have a quantifiable
benefit in predicting RUL on the NASA CMAPSS dataset.

\section{Methods}

\subsection{Dataset and Preprocessing}

We utilize a benchmark dataset named NASA CMAPSS FD001. This dataset
is a multivariate time series dataset that contains information collected from
multiple jet engines running under a single condition with one fault mode. Each
row in this dataset corresponds to a single engine cycle. It contains information
about an engine's identifier, cycle number, three operational setting values, and
twenty-one sensor values.

The training dataset has complete run-to-failure trajectories for every single
engine. This enables us to calculate the remaining useful life. In this dataset,
we calculate the remaining useful life for every single time step. This is done
by subtracting the current cycle number from the maximum cycle number for
every single engine.

In the pre-processing step, we loaded the raw data files. This is done using
a whitespace-separated parsing mechanism. Trailing empty columns were also
removed. Every single column is given a specific name corresponding to every
single engine's identifier, cycle number, operational setting values, and sensor
values.

\subsection{Feature Engineering}

To build the input features, the engine ID and cycle number were removed to
prevent the use of identification variables or absolute time as a predictor. The
remaining operational variables and sensor readings were used as the features.
Standardization of the numeric features was performed by applying a standard
scaling technique.

\subsection{Models}

Two baseline regression models were implemented for comparison:

\begin{itemize}
\item Linear Regression: A simple model that assumes a linear relationship
between input features and the target variable.
\item Random Forest Regression: An ensemble-based model that captures
nonlinear relationships and interactions between features.
\end{itemize}

To prevent data leakage across time steps coming from the same engine, an
engine-level train validation split is used. This ensures that all cycles from a
given engine occur within only one split.

\subsection{Evaluation Metrics}

Model performance was evaluated using two standard regression metrics:

\begin{itemize}
\item Mean Absolute Error (MAE): Measures the average absolute differ-
ence between predicted and true RUL values.
\item Root Mean Squared Error (RMSE): Measures the square root of the
average squared differences, placing greater emphasis on larger errors.
\end{itemize}

\section{Preliminary Results}

Preliminary experiments were performed on the NASA CMAPSS FD001 dataset
to assess baseline performance on the Remaining Useful Life (RUL) prediction
task. This dataset has 20,631 training samples and 13,096 test samples, each of
which is associated with 100 engines.

Two baseline models, Linear Regression and Random Forest Regression,
were trained on standardized features from sensors and operations. For evalua-
tion, a validation split was used on an engine level to prevent leakage.

The results are given below:

\begin{itemize}
\item Random Forest Regression: MAE = 25.874, RMSE = 35.400
\item Linear Regression: MAE = 30.111, RMSE = 38.133
\end{itemize}

The results show the superiority of the Random Forest model over Linear
Regression, which confirms the importance of nonlinear relationships in the
sensor data for the prediction of RUL. However, the results show the limitation
of both models in handling the temporal degradation patterns in the data.

The above results show the importance of using sequence-based models, such
as Convolutional Neural Networks (CNNs), which are more efficient in handling
the temporal dependencies in the engine degradation data. Further work will
be done in this direction to improve the prediction accuracy.

\section{Future plan}

On the basis of these preliminary results, further research will be conducted to
enhance the performance of the models by incorporating techniques for tempo-
ral modeling. For this purpose, sequence-based models, such as Convolutional
Neural Networks (CNN), will be incorporated for better degradation pattern
recognition. Further improvements in this research will be made by trying out
different window sizes for sequence generation, tuning the parameters of the
models, and exploring feature selection techniques for identifying the best sen-
sors. Further evaluation of the models will be conducted by employing both
MAE and RMSE for assessing the models' performance against different error
measures. Moreover, a comparative study of baseline models and deep learning
models will be conducted to evaluate the effectiveness of temporal models in
enhancing the accuracy of predicted Remaining Useful Life.

\section{References}

References

[1] A. Saxena, K. Goebel, D. Simon, and N. Eklund, ``Damage Propaga-
tion Modeling for Aircraft Engine Run-to-Failure Simulation,'' International
Conference on Prognostics and Health Management, 2008.

[2] S. Zheng, K. Ristovski, A. Farahat, and C. Gupta, ``Long Short-Term Mem-
ory Network for Remaining Useful Life Estimation,'' IEEE International
Conference on Prognostics and Health Management, 2017.

[3] X. Li, Q. Ding, and J. Sun, ``Remaining Useful Life Estimation in Prognos-
tics Using Deep Convolution Neural Networks,'' Reliability Engineering \&
System Safety, 2018.


\end{document}
